\noindent\markforreview
\vspace{1em}

Given the expression
\[
\sqrt[5]{70n \left(\sqrt[7]{70n}\right)^2}
\]
for what value of \(x\) is the given expression equivalent to \((70n)^{30x}\), where \(n > 1\)?

\vspace{2em}
\newpage

\noindent\markforreview
\vspace{1em}

The expression \(4x^2 + bx - 45\), where \(b\) is a constant, can be rewritten as
\[
(hx + k)(x + j)
\]
where \(h, k,\) and \(j\) are integer constants.

Which of the following must be an integer?

\vspace{0.75em}

\choicebox{A}{\(\dfrac{b}{h}\)}
\choicebox{B}{\(\dfrac{b}{k}\)}
\choicebox{C}{\(\dfrac{45}{h}\)}
\choicebox{D}{\(\dfrac{45}{k}\)}

\vspace{2em}
\newpage

\noindent\markforreview
\vspace{1em}

The expression
\[
6\sqrt[5]{3x^{45}} \cdot \sqrt[8]{8x}
\]
is equivalent to \(ax^b\), where \(a\) and \(b\) are positive constants and \(x > 1\).

What is the value of \(b\)?

\vspace{2em}
\newpage

\noindent\markforreview
\vspace{1em}

In the given equation, \(c\) is a positive constant. Solve for \(x\):
\[
\frac{x^2}{\sqrt{x^2 - c^2}} = \frac{c^2}{\sqrt{x^2 - c^2}} + 39
\]

Which of the following is one of the solutions to the given equation?

\vspace{0.75em}

\choicebox{A}{\(-c\)}
\choicebox{B}{\(-c^2 - 39^2\)}
\choicebox{C}{\(-\sqrt{39^2 - c^2}\)}
\choicebox{D}{\(\sqrt{c^2 + 39^2}\)}

\vspace{2em}
\newpage

\noindent\markforreview
\vspace{1em}

The quadratic function is given as
\[
f(x) = ax^2 + 4x + c
\]
where \(a\) and \(c\) are constants. The graph of \(y = f(x)\) is a parabola that opens upward and has a vertex at \((h, k)\), where \(h\) and \(k\) are constants.

If \(k < 0\) and \(f(-9) = f(3)\), which of the following must be true?

\begin{itemize}
\item[I] \(c < 0\)
\item[II] \(a \geq 1\)
\end{itemize}

\vspace{0.75em}

\choicebox{A}{I only}
\choicebox{B}{II only}
\choicebox{C}{I and II}
\choicebox{D}{Neither I nor II}

\vspace{2em}
\newpage

\noindent\markforreview
\vspace{1em}

The functions \(f\) and \(g\) are defined by the given equations, where \(x \geq 0\):

\begin{itemize}
\item[I] \(f(x) = 33(0.4)^{x+3}\)
\item[II] \(g(x) = 33(0.16)(0.4)^{x-2}\)
\end{itemize}

Which of the following equations displays, as a constant or coefficient, the maximum value of the function it defines, where \(x \geq 0\)?

\vspace{0.75em}

\choicebox{A}{I only}
\choicebox{B}{II only}
\choicebox{C}{I and II}
\choicebox{D}{Neither I nor II}

\vspace{2em}
\newpage

\noindent\markforreview
\vspace{1em}

Each of the following frequency tables represents a data set.

Which of these frequency tables represents the data set with the \textbf{smallest standard deviation}?

\[
\begin{array}{|c|c|c|c|c|}
\hline
\text{Value} & \text{Frequency (A)} & \text{Frequency (B)} & \text{Frequency (C)} & \text{Frequency (D)} \\
\hline
5 & 0 & 11 & 15 & 18 \\
6 & 10 & 17 & 15 & 14 \\
7 & 55 & 19 & 15 & 11 \\
8 & 10 & 17 & 15 & 14 \\
9 & 0 & 11 & 15 & 18 \\
\hline
\end{array}
\]

\vspace{0.75em}

\choicebox{A}{A}
\choicebox{B}{B}
\choicebox{C}{C}
\choicebox{D}{D}

\vspace{2em}
\newpage

\noindent\markforreview
\vspace{1em}

In \(\triangle ABC\), \(AB = 4\) and \(BC = 5\). If \(\angle ABC < 90^\circ\), which of the following could be the area of \(\triangle ABC\)?

\vspace{0.75em}

\choicebox{A}{5}
\choicebox{B}{10}
\choicebox{C}{15}
\choicebox{D}{20}

\vspace{2em}
\newpage

\noindent\markforreview
\vspace{1em}

The table summarizes the distribution of heights, in inches (\("in"\)), for 300 participants in a study.

If a participant that is at most 71 inches tall is chosen at random, what is the probability that the participant was selected from Group 2?

\[
\begin{array}{|c|c|c|c|c|}
\hline
 & \text{48--59 in.} & \text{60--71 in.} & \text{72--83 in.} & \text{Total} \\
\hline
\text{Group 1} & 20 & 60 & 20 & 100 \\
\text{Group 2} & 10 & 75 & 15 & 100 \\
\text{Group 3} & 30 & 45 & 25 & 100 \\
\hline
\text{Total} & 60 & 180 & 60 & 300 \\
\hline
\end{array}
\]

\vspace{0.75em}

\choicebox{A}{\(\dfrac{1}{4}\)}
\choicebox{B}{\(\dfrac{17}{48}\)}
\choicebox{C}{\(\dfrac{17}{60}\)}
\choicebox{D}{\(\dfrac{17}{20}\)}

\vspace{2em}
\newpage

\noindent\markforreview
\vspace{1em}

The function \(f\) is defined as \(f(x) = x^2 + ax + b\), where \(a\) and \(b\) are positive integers.

If the minimum value of \(f(x)\) is 7, which is a possible value of \(b\)?

\vspace{0.75em}

\choicebox{A}{6}
\choicebox{B}{7}
\choicebox{C}{8}
\choicebox{D}{9}

\vspace{2em}
\newpage

\noindent\markforreview
\vspace{1em}

The average annual energy cost for a certain home is \$4{,}334. The homeowner plans to spend \$25{,}000 to install a geothermal heating system. The homeowner estimates that the average annual energy cost will then be \$2{,}712.

Which of the following inequalities can be solved to find \(t\), the number of years after installation at which the total amount of energy cost savings will exceed the installation cost?

\vspace{0.75em}

\choicebox{A}{\(25{,}000 > (4{,}334 - 2{,}712)t\)}
\choicebox{B}{\(25{,}000 < (4{,}334 - 2{,}712)t\)}
\choicebox{C}{\(25{,}000 - 4{,}334 - 2{,}712t\)}
\choicebox{D}{\(25{,}000 > \dfrac{4{,}332}{2{,}712}t\)}

\vspace{2em}
\newpage

\noindent\markforreview
\vspace{1em}

For the equation, \(b\) and \(c\) are constants. For any real number \(t\), the point
\[
\left(\frac{-2t + 4}{3},\, t\right)
\]
lies on the graph of the equation \(bx + 3y = c\).

What is the value of \(c\)?

\vspace{0.75em}

\choicebox{A}{-2}
\choicebox{B}{4}
\choicebox{C}{6}
\choicebox{D}{12}

\vspace{2em}
\newpage

\noindent\markforreview
\vspace{1em}

A circle in the \(xy\)-plane has equation
\[
(x - 2.5)^2 + (y - 3.5)^2 = 225.
\]

Line \(p\) is tangent to the circle at point \((-6.5, C)\), where \(C\) is a constant.

Which of the following could be the slope of line \(p\)?

\begin{itemize}
\item[I] \(-\dfrac{3}{4}\)
\item[II] \(\dfrac{3}{4}\)
\item[III] \(\dfrac{4}{3}\)
\end{itemize}

\vspace{0.75em}

\choicebox{A}{I only}
\choicebox{B}{I and II only}
\choicebox{C}{I and III only}
\choicebox{D}{I, II, and III}

\vspace{2em}
\newpage

\noindent\markforreview
\vspace{1em}

The function \(f\) is defined by \(f(x) = a^x - b\), where \(a\) and \(b\) are constants.

In the \(xy\)-plane, the graph of \(y = f(x)\) passes through the points \((c, 7)\) and \((2c, 117)\), where \(c\) is a constant.

What is the value of \(b\)?

\vspace{2em}
\newpage

\noindent\markforreview
\vspace{1em}

The given equation, where \(t\) is a positive constant, defines a circle in the \(xy\)-plane.

The radius of this circle is \(\sqrt{41}\).

What is the value of \(t\)?
\[
2x^2 + 2y^2 + tx + ty - 78 = 0
\]
